\section{Exact implementation and complete verifiers}\label{app:verifiers}

This appendix identifies the array operations with the certificates of
Section~\ref{sec:polynomial-certificates} and bounds the auxiliary
fixed-width arithmetic.  The complete source listings follow these
implementation details.  The command lines and expected C++ output are
given in Section~\ref{sec:finite-verification}.

\subsection{Exact coefficient arithmetic}\label{app:implementation}

The C++ type \texttt{P} represents a rational coefficient matrix of bidegree
$(r,r)$ by an integer array and a common positive denominator.  Its entries
are indexed by $\mathcal I_r$ in descending lexicographic order.  Integer
numerators and denominators use \texttt{boost::multiprecision::cpp\_int}.
The routine \texttt{ones(r)} constructs $W_r$; multiplication is ordinary
coefficient convolution, and \texttt{lift} raises both degrees by multiplying
by $W_1$.  Addition first lifts to a common degree and then uses a common
denominator.  Thus signs are tested on integer numerators, without a
floating-point tolerance.

We explain the correspondence between the remaining routines and the
mathematical certificates; this also fixes the normalizations in the
listing.  Initially the coordinates of every vertex are stored as integers
at scale
\begin{equation}\label{eq:vertex-scale}
 \kappa=12\cdot2^{24}.
\end{equation}
Removing the factor $2^{24}$ leaves the denominator 12 for the individual
coordinates.  The coefficients of $U_v$ are
$4(a_i)_v(b_j)_{\pi(v)}$; their initialization removes $2^{48}$ and uses
denominator 144.  The routine \texttt{inv} forms the nine invariant
polynomials in their stated order.  The routine \texttt{dot} forms a linear
combination divided by 1000 and homogenizes lower-degree terms with its
third argument.  With that argument equal to the homogeneous one it gives
$H(p)$, $F(p)$, or $G(p)$; with that argument equal to $z$ it gives
$z^6H(U/z)$.

There is one useful economy in constructing $K$.  A polynomial depending
only on $X$ is nevertheless stored with both degrees equal, so it has the
form $f(X)(\sum_jY_j)^r$.  Its column indexed by $(r,0,0,0)$ contains exactly
the coefficients of $f$.  Taking this column and the analogous row of a
polynomial depending only on $Y$ gives the outer product used by
\texttt{tensor}.  The factor 24 in \texttt{make} then gives exactly the $K$
of \eqref{eq:rational-K}.  The array \texttt{t} is the Type T polynomial;
it is constructed once and thereafter restricted by midpoint substitution.
The arrays \texttt{N} and \texttt{D} store $N_\pi$ and $z_\pi$, respectively,
for the 24 permutations in lexicographic order.

The Type P interpolation uses ordinary matrix multiplication.  The matrix
\texttt{ev(r,3)} has entries $I^J$, where $I\in\mathcal I_3$ indexes a
node and $J\in\mathcal I_r$ indexes a monomial.  Consequently the two-sided
evaluation \texttt{act} must be divided by $3^{2r}$.  The listing includes
this factor explicitly, using 9 when $r=1$.  The routine \texttt{inverse}
constructs, rather than numerically inverts, the cardinal polynomials
in \eqref{eq:interpolation-polynomial}.  It expands their homogeneous
linear factors $3X_h-m\sum_iX_i$, stores their coefficients multiplied by
960, and divides by the factorial denominators.  Two-sided interpolation
therefore requires division by $960^2$.

After the node values of $N_\pi/z_\pi^k$ have been truncated toward zero in
multiples of $M^{-1}$, \texttt{testP} constructs the interpolant $J_\pi$.
Its subtraction \texttt{o=b-o} implicitly lifts $N_\pi$ by $W_2$, and hence
constructs precisely $O_\pi=W_2N_\pi-z_\pi^kJ_\pi$.
The array \texttt{w} is $Q_\pi=z_\pi^kW_3$.  For each negative coefficient
of $O_\pi$, the program asserts that the corresponding coefficient of
$Q_\pi$ is positive.  A violation aborts the verification rather than
merely rejecting that cell; no such violation occurs in the completed
runs.  It then computes the integer ceiling in
\eqref{eq:remainder-correction}, by positive integer division, and retains
the largest value.  The final test is exactly the coefficient test of
Lemma~\ref{lem:type-p}.

Finally, \texttt{sub} implements \eqref{eq:midpoint-substitution} with
$2^r$ cleared.  Its integer multiplier is
\begin{equation}\label{eq:cleared-substitution}
 C_t=2^{r-I_i-t}\binom{I_j}{t},\qquad
 C_{t+1}=\frac{C_t(I_j-t)}{2(t+1)}.
\end{equation}
Here $r-I_i-t\ge0$ for every used term, so the indicated division is exact.
The resulting matrix denominator is multiplied by $2^r$.

\subsection{Bounds for fixed-width arithmetic}

Only indices, counters, vertex coordinates, and small combinatorial factors use
fixed-width integer types.  Take \texttt{int} to have at least 32 bits and
\texttt{long long} to have at least 64 bits.  The following bounds cover
every such operation in the listing.

All polynomial matrices in the C++ verifier have bidegree at most 11.
Thus there are at most $\binom{14}{3}=364$ indices in one dimension and
$132\,496$ entries in a matrix.  The depth limit 24 allows at most
$2^{25}-1$ visited cells, so the counters for visited cells and terminal
cells also fit a 32-bit signed integer.  The largest factorial initialized is
$19!=121\,645\,100\,408\,832\,000<2^{63}$, and products of factorials
appearing in a multinomial denominator do not exceed this bound.  The
largest multinomial coefficient through degree 11 is $92\,400$, whose
square is $8\,537\,760\,000<2^{63}$.

Each scaled vertex coordinate belongs to $[0,\kappa]$.  Initial products
are bounded by $4\kappa^2$, and even the conservative bound
\[
 64\kappa^2
 =2\,594\,073\,385\,365\,405\,696<2^{63}
\]
covers the weighted squared edge distances.  The initial coordinates are
multiples of $2^{24}$, so all midpoint coordinates are integers through
depth 24.  A point of $\mathcal S$ has fourth coordinate at most $1/4$, with
equality only at $\unif$.  Thus the test that a stored fourth coordinate is
$\kappa/4$ detects a uniform vertex exactly.

The multipliers in \eqref{eq:cleared-substitution}, including the
intermediate numerator of the recurrence, are bounded by
$11\cdot2^{11}\binom{11}{5}<2^{24}$.  The evaluation matrices contain
integers at most $3^6$, and the cardinal-polynomial construction uses only
three linear factors starting from 960.  Even bounding the sum of absolute
coefficients at each multiplication by a factor 8 gives
$960\cdot8^3<2^{19}$.  All products with polynomial coefficients are already
operations on unlimited integers.  These bounds exclude fixed-width
overflow; in particular, the exact sign tests do not rely on floating-point
arithmetic or machine rounding.

\subsection{Rational sign certificates}
\begingroup
\interlinepenalty=10000
\lstinputlisting[language=C++,basicstyle=\ttfamily\scriptsize,
  columns=fullflexible,keepspaces=true,breaklines=true,
  showstringspaces=false,numbers=left,numberstyle=\tiny,
  numbersep=6pt,tabsize=4]{../verification/checks/certificate.cpp}
\endgroup

\subsection{Scalar coefficients and symmetry}
\lstinputlisting[language=Python,basicstyle=\ttfamily\scriptsize,
  columns=fullflexible,keepspaces=true,breaklines=true,
  showstringspaces=false,numbers=left,numberstyle=\tiny,
  numbersep=6pt,tabsize=4]{../verification/checks/small_polynomials.py}
